\documentclass{standalone}
\usepackage{tikz}
\usetikzlibrary{calc,through}
% i have no ideo how this works
\tikzset{circumcircle/.style n args={3}{%
insert path={let    \p1=($(#1)!0.5!(#2)$),
                    \p2=($(#1)!0.5!(#3)$),
                    \p3=($(#1)!0.5!(#2)!1!-90:(#2)$),
                    \p4=($(#1)!0.5!(#3)!1!90:(#3)$),
                    \p5=(intersection of \p1--\p3 and \p2--\p4)
                    in },
at={(\p5)},
circle through= {(#1)},circle={0.1cm}
}}
\def\scalep{0.4}
\begin{document}
    \begin{tikzpicture}
    % clip a tikzpicture
        \clip (-10*\scalep,-0.5) rectangle (10*\scalep,100*\scalep*\scalep);
        \draw[domain = -5*\scalep:10*\scalep, variable = \x] plot ({\x}, {\x*\x*0.5});
        % no spaces in foreach, otherwise, they are included in the variables
        \foreach \x/\nodename in {1/A,2/B,3/C,4/D,5/E,6/F,7/G,8/H,9/I}
        {
            \ifodd\x
                \draw[fill, color=red!50!red] (\x*\scalep, \x*\x*\scalep*\scalep*0.5) circle (\scalep*0.4cm) node (\nodename){};
            \else
                \draw[fill, color=blue!50!blue] ({\x*\scalep}, {\x * \x*\scalep*\scalep*0.5}) circle (\scalep*0.4cm) node (\nodename){};
            \fi
        }
        \node[circumcircle={B}{F}{H}, draw=red] {};
    \end{tikzpicture}
\end{document}
